\documentclass[11pt,a4paper]{article}

% ========================
% ESSENTIAL PACKAGES
% ========================

% Language and encoding
\usepackage[utf8]{inputenc}
\usepackage[T1]{fontenc}
\usepackage[english]{babel}
%\usepackage{natbib}

% Math packages
\usepackage{amsmath}        % Enhanced math environments
\usepackage{amssymb}        % Additional math symbols
\usepackage{amsthm}         % Theorem environments
\usepackage{amsfonts}       % Additional math fonts
\usepackage{mathtools}      % Extension to amsmath
\usepackage{bm}             % Bold math symbols
\usepackage{dsfont}         % Double-struck fonts (for sets like ℝ, ℕ)
\usepackage{mathrsfs}       % Script fonts for math

% Graphics and figures
\usepackage{graphicx}       % Include graphics
\usepackage{float}          % Better float placement
\usepackage{subcaption}     % Subfigures
\usepackage{tikz}           % Drawing diagrams
\usepackage{pgfplots}       % Plotting
\pgfplotsset{compat=1.18}

% Page layout and formatting
\usepackage[margin=1in]{geometry}
\usepackage{fancyhdr}       % Custom headers and footers
\usepackage{titlesec}       % Customize section titles
\usepackage{enumitem}       % Customize lists
\usepackage{parskip}        % Paragraph spacing

% Colors and highlighting
\usepackage{xcolor}
\usepackage{mdframed}       % Framed environments
\usepackage{tcolorbox}      % Colored boxes

% Tables
\usepackage{array}          % Extended array environments
\usepackage{booktabs}       % Professional tables
\usepackage{multirow}       % Multi-row cells

% References and links
\usepackage{hyperref}       % Hyperlinks
\usepackage{cleveref}       % Smart cross-references

% Additional useful packages
\usepackage{cancel}         % Cancel terms in equations
\usepackage{siunitx}        % SI units
\usepackage{listings}       % Code listings (if needed)

% ========================
% CUSTOM COMMANDS
% ========================

% Common math sets
\newcommand{\R}{\mathbb{R}}
\newcommand{\N}{\mathbb{N}}
\newcommand{\Z}{\mathbb{Z}}
\newcommand{\Q}{\mathbb{Q}}
\newcommand{\C}{\mathbb{C}}

% Common operators
\DeclareMathOperator{\lcm}{lcm}
\DeclareMathOperator{\Tr}{Tr}
\DeclareMathOperator{\rank}{rank}
%\DeclareMathOperator{\span}{span}
%\DeclareMathOperator{\null}{null}

% Shortcuts for common expressions
\newcommand{\abs}[1]{\left|#1\right|}
\newcommand{\norm}[1]{\left\|#1\right\|}
\newcommand{\inner}[2]{\left\langle #1, #2 \right\rangle}
\newcommand{\ceil}[1]{\left\lceil #1 \right\rceil}
\newcommand{\floor}[1]{\left\lfloor #1 \right\rfloor}

% Derivatives and integrals
\newcommand{\dd}[1]{\,\mathrm{d}#1}
\newcommand{\dv}[2]{\frac{\mathrm{d}#1}{\mathrm{d}#2}}
\newcommand{\pdv}[2]{\frac{\partial #1}{\partial #2}}

\renewcommand{\hat}{\widehat}
\renewcommand{\tilde}{\widetilde}

% ========================
% THEOREM ENVIRONMENTS
% ========================

\theoremstyle{definition}
\newtheorem{definition}{Definition}[section]
\newtheorem{example}{Example}[section]
\newtheorem{exercise}{Exercise}[section]

\theoremstyle{plain}
\newtheorem{theorem}{Theorem}[section]
\newtheorem{lemma}{Lemma}[section]
\newtheorem{corollary}{Corollary}[section]
\newtheorem{proposition}{Proposition}[section]
\newtheorem{assumption}{Assumptio}[section]

\theoremstyle{remark}
\newtheorem{remark}{Remark}[section]
\newtheorem{note}{Note}[section]

\newcommand{\supp}{\mathrm{supp}}
\newcommand{\op}{\mathrm{op}}

\newcommand{\mX}{\mathcal X}
\newcommand{\mD}{\mathcal D}
\newcommand{\bP}{\mathbb P}
\newcommand{\lp}{L^2(\bP)}
\newcommand{\lx}{L^2(\mX)}
\newcommand{\ltwo}{L^2(\mX)}
\newcommand{\ud}{\mathrm d}
\newcommand{\mH}{\mathcal H}
\newcommand{\mA}{\mathcal A}

\newcommand{\diag}{\mathrm{diag}}

% ========================
% COLORED BOXES
% ========================

% Important theorem box
\newtcolorbox{importantbox}{
	colback=blue!5!white,
	colframe=blue!75!black,
	title=Important
}

% Warning/Note box
\newtcolorbox{notebox}{
	colback=yellow!10!white,
	colframe=orange!75!black,
	title=Note
}

% Example box
\newtcolorbox{examplebox}{
	colback=green!5!white,
	colframe=green!75!black,
	title=Example
}

% ========================
% DOCUMENT SETUP
% ========================

% Page style
\pagestyle{fancy}
\fancyhf{}
\rhead{\thepage}
\lhead{\leftmark}

% Title formatting
\title{Functional Bandit Notes}
\author{Your Name}
\date{\today}

% Hyperlink setup
\hypersetup{
	colorlinks=true,
	linkcolor=blue,
	filecolor=magenta,      
	urlcolor=cyan,
	citecolor=red
}

\begin{document}

The main purpose of this document is to give a technical summary of existing approaches towards (offline) functional regression,
and remark on their strengths/weakness. When appropriate, I will also remark on possibilities to bandit problems (mostly confidence intervals).

Let $\mX$ be a compact Euclidean set and $\mD=L^2(\mX)$ be the set of all square integrable functions on $\mX$.
Let $\mathbb P$ be the distribution on $\mD$. The model is 
$$
y_i = \langle f_i, \beta_0\rangle + \varepsilon_i
$$
where $\{\varepsilon_i\}_{i=1}^n$ are centered, exogenous sub-Gaussian random variables, $\beta_0\in\mathcal F\subseteq \mD$ is an unknown regression model and $f_1,\cdots,f_N\overset{i.i.d.}{\sim}\mathbb P$.
The inner product is taken with respect to the Lesbesgue measure on $\mX$:
$$
\langle f_i, \beta_0\rangle := \int_{\mX} f_i(u)\beta_0(u)\ud u.
$$
Sometimes, we will place additional assumptions on $\mathcal F$, or restrict the support of $\bP$.
Errors are measured in 
$$
\|\hat \beta_n-\beta_0\|_{\lp}^2 := \int_{\mD} \langle f, \hat \beta_n-\beta_0\rangle^2\ud\bP(f).
$$

Let $K$ be the RKHS and recall the notation that $\tilde f = L_{K^{1/2}}f$, $\tilde\beta_0=L_{K^{-1/2}}\beta_0$, such that $\|\tilde\beta_0\|_2 < +\infty$.
Let $T=\mathbb E[\tilde f\otimes \tilde f]$ be the population covariance operator, which is decomposed as $T=\sum_{j=1}^{\infty}s_j\psi_j\otimes\psi_j$ for an orthonormal basis $\{\psi_j\}_{j=1}^{\infty}$ of $L_2(\mX)$,
and $s_j\asymp j^{-2r}$.

\subsection{The anisotropic KRR approach}

Given observations $\{f_i,y_i\}_{i=1}^n$, let 
$$
T_n := \frac{1}{n}\sum_{i=1}^n \tilde f_i\otimes\tilde f_i
$$
be the sample covariance operator, where $\tilde f_i = L_{K^{1/2}}f_i$. By the spectral theorem, there exist orthonormal (w.r.t.~$L^2(\mX)$) $\{\hat\psi_j\}_{j=1}^n$ and $\hat s_1\geq\cdots\hat s_n\geq 0$ such that
$$
T_n = \sum_{j=1}^n \hat s_j\hat\psi_j\otimes\hat\psi_j.
$$
Note that $T_n$ can only have at most $n$ non-zero eigenvalues.

Given a function $f$, let $\hat\alpha_j := \langle \tilde f, \hat\psi_j\rangle=\langle L_{K^{1/2}}f,\hat\psi_j\rangle$. We then use the following anisotropic KRR estimator
$$
\hat\beta = \left(T_n + \Lambda_n(f)\right)^{-1}\cdot \frac{1}{n}\sum_{i=1}^n y_i L_{K^{1/2}}f_i,
$$
where
$$
\Lambda_n(f) = \sum_{j=1}^n\lambda_j(f)\hat\psi_j\otimes\hat\psi_j;\;\;\;\;\;\;\lambda_j(f) = \left\{\begin{array}{ll} \hat s_j\sqrt{\lambda_0}/|\hat\alpha_j|,& j\leq J_1;\\ +\infty,& j>J_1,\end{array}\right.
$$
where $J_1\leq n$ is a cut-off threshold to be determined later.
{\color{red}Note that this is a truncation + KRR approach.}
For notational convenience, we will drop the dependency on $f$ whenever there is no ambiguity. An estimate of $\langle f,\beta_0\rangle$ is the obtained as $\langle L_{K^{1/2}} f, \hat\beta\rangle$.

\subsection{Analysis of bias}

Recall the definition that $\tilde\beta_0=L_{K^{-1/2}}\beta_0$. Let 
$$
\bar\beta := \mathbb E[\hat\beta] = (T_n+\Lambda_n)^{-1}\frac{1}{n}\sum_{i=1}^n \langle\tilde f_i,\tilde\beta_0\rangle \tilde f_i = (T_n+\Lambda_n)^{-1}T_n\tilde\beta_0.
$$
Let $\hat\theta_j := \langle\tilde\beta_0, \hat\psi_j\rangle$. Then
$$
\tilde\beta_0-\bar\beta = (T_n+\Lambda_n)^{-1}\Lambda_n \tilde\beta_0 = \sum_{j=1}^n \frac{\lambda_j\hat\theta_j}{\hat s_j+\lambda_j}\hat\psi_j.
$$
Subsequently, the bias is upper bounded by
$$
\big|\langle \tilde f,\tilde\beta_0-\bar\beta\rangle\big| = \left|\sum_{j=1}^n\frac{\lambda_j\hat\theta_j\hat\alpha_j}{\hat s_j+\lambda_j}\right|.
$$

From the definition of $\lambda_j$, it holds for every $j\in[n]$ that
$$
\frac{\lambda_j|\hat\theta_j\hat\alpha_j|}{\hat s_j+\lambda_j} \leq \left\{\begin{array}{ll} \sqrt{\lambda_0}|\hat\theta_j|,& j\leq J_1;\\ |\hat\theta_j\hat\alpha_j|,& j>J_1.\end{array}\right.
$$
Let $J_1 := n^{4r/(4r^2-1)} = J_0^{2r/(2r-1)}$ where $J_0=n^{1/(2r+1)}=\lambda_0^{-1/(2r)}$. We use the following approach to upper bound the RHS of the above inequality:
\begin{align*}
\left|\sum_{j=1}^n\frac{\lambda_j\hat\theta_j\hat\alpha_j}{\hat s_j+\lambda_j}\right|
&\leq \sqrt{\lambda_0}\left(\sum_{j=1}^{J_1}|\hat\theta_j|\right) + \sqrt{\sum_{j=J_1+1}^n \hat\theta_j^2}\sqrt{\sum_{j=J_1+1}^n|\hat\alpha_j|^2}\\
&\leq \sqrt{\lambda_0}\|\hat\theta_{J_1}\|_1 + \lambda_0\|\hat\theta\|_2\|\hat\alpha_{(J_1+1):n}\|_2.
\end{align*}
It should also be noted that, if $\sum_{j=J_1+1}^n \hat s_j\lesssim \sum_{j=J_1+1}^n s_j$, then
$$
\mathbb E[\|\hat\alpha_{(J_1+1):n}\|_2^2] \lesssim \int_{J_1}^{\infty} t^{-2r}\ud t \lesssim J_1^{1-2r} = J_0^{-2r} = \lambda_0.
$$

\subsection{Analysis of perturbation}

We impose the following two main assumptions, on ``population'' coefficients $\theta_j := \langle \tilde\beta_0,\psi_j\rangle$:
$$
\sum_{j=1}^{\infty}\theta_j^2 < +\infty,\;\;\;\;\;\;\sum_{j=1}^{\infty}|\theta_j| < +\infty.
$$
The first condition is standard which simply means $\|\tilde\beta_0\|_2<+\infty$. This condition is also basis-invariance so applicable to $\{\hat\theta_j\}_j$ too. The main issue is the second assumption, which is \emph{not} basis invariant.

We state the following result, which controls the behaviors of $\{\hat\theta_j\}$ and $\{\hat s_j\}$.
\begin{lemma}
Suppose the following additional conditions hold for every $j\in[n]$:
\begin{enumerate}
\item $s_j\asymp j^{-2r}$ and $s_j-s_{j+1}\gtrsim j^{-2r-1}$;
\item The random variable $\hat\alpha_j = \langle \tilde f,\hat\psi_j\rangle$, with randomness on $\tilde f$ only, is sub-Gaussian with variance proxy $s_j$.
\end{enumerate}
Then, if $r>1+\sqrt{5}/2=2.118=:r_0$, the following holds with probability $1-o(1)$:
\begin{enumerate}
\item $\sum_{j=1}^{J_1}|\hat\theta_j| \leq \sum_{j=1}^{J_1}|\theta_j| + o(1)$;
\item $\hat s_j\gtrsim j^{-2r}$ for all $1\leq j\leq J_1$;
\item $\sum_{j=J_1+1}^n \hat s_j \lesssim J_1^{-2r+1}$.
\end{enumerate}
\label{lem:hat-psi}
\end{lemma}

\begin{remark}
The sub-Gaussianity assumption and the spectral gap assumption can be somewhat relaxed, at the cost of a larger value of $s_0$. It should be noted that both assumptions are adopted for fPCA. We do \emph{not} adopt the fPCA assumption of decays on $\{\theta_j\}$, however, because that leads to significant sub-optimality.
\end{remark}

\begin{proof}[Proof of Lemma \ref{lem:hat-psi}.]
Express both $T$ and $T_n$ on the first $n$ eigenfunctions $\{\psi_i\}_{i=1}^n$ of $T$. Then $T=\diag(s_1,\cdots,s_n)$ with the standard basis as its eigenvectors, and $T_n$ is a PSD matrix. 
We will carry out our proof in this $n\times n$ finite matrix field to simplify matters.

We will use the following inequality as the key workhorse of this proof: with probability $1-o(1)$ uniformly for every $i,j\in[n]$, 
$$
\big|[T_n]_{ij} - T_{ij}\big| \lesssim \frac{i^{-r}j^{-r}}{\sqrt{n}}.
$$
We can potentially tolerate higher-order terms in this inequality, if we want to relax the sub-Gaussianity assumption, but Bernstein with the simple $|T_{ij}|\leq 1$ is \emph{not} enough.
The rest of the proof are conditioned on the success event that the above holds.

Consider the submatrix of $T_n$ corresponding to $\psi_{J_1+1},\cdots,\psi_n$. The trace formula implies that
\begin{align*}
\sum_{j=J_1+1}^{n} \hat s_j &\leq \sum_{j=J_1+1}^n [T_n]_{jj} \lesssim \sum_{j=J_1+1}^n j^{-2r} + O\left(\frac{j^{-2r}}{\sqrt{n}}\right) \lesssim J_1^{-2r+1},
\end{align*}
which proves the third property.

Now, let $A$ and $A_n$ be the $J_1\times J_1$ sub-matrices of $T$ and $T_n$ corresponding to the eigenspace spanned by $(\psi_1,\cdots,\psi_{J_1})$.
We will prove that, for every $j\in[J]$, there exists $w\in\mathbb R^d$ satisfying $\|w\|_2 = o(J_1^{-1})$ such that
\begin{equation}
A_n (e_j+w) = s_j(e_j+w).
\label{eq:eigen-an}
\end{equation}
This shows that $\hat s_j=s_j(1+o(J_1^{-1}))$. Because the eigengaps of the first $J_1$ eigenfunctions are larger than $J_1^{-1}$, this ensures that the top-$J_1$ eigenfunctions of $T$ and $T_n$ have the same span $\mathrm{span}(\psi_1,\cdots,\psi_{J_1})$, and furthermore $\|\hat\psi_j-\psi_j\|_2 = o(J_1^{-1})$. This is sufficient to prove the first two properties.

Let $H=A_n-A$ and consider $w_j=0$ which is simply a Gauge choice. Because $Ae_j=s_je_j$, Eq.~(\ref{eq:eigen-an}) implies
$$
(A-s_j I+H)w = -He_j.
$$
Let $\bar H\in\mathbb R^{(J_1-1)\times (J_1-1)}$ be the submatrix formed by removing the $j$th row and $j$th column of $H$. Let $h\in\mathbb R^{J_1-1}$ be the $j$th column of $H$, removing $H_{jj}$.
Let $D=\diag(s_1-s_j,\cdots,s_{j-1}-s_j,s_{j+1}-s_j,\cdots,s_{J_1}-s_j)\in\mathbb R^{(J_1-1)\times(J_1-1)}$. Let $\bar w=(w_1,\cdots,w_{j-1},w_{j+1},\cdots,w_n)\in\mathbb R^{J_1-1}$. Then
$$
\bar w = -(D+\bar H)^{-1}h = -D^{-1/2}(I+\tilde H)^{-1}D^{-1/2}h,
$$
where $\tilde H = D^{-1/2} \bar H D^{-1/2}$. 

For any $i,k\neq j$, we have that
$$
\big|\tilde H_{ik}\big| \lesssim \frac{1}{\sqrt{n}} \frac{i^{-r}k^{-r}}{\sqrt{|(s_i-s_j)(s_k-s_j)|}} =: \frac{1}{\sqrt{n}}z_iz_k\;\;\;\;\;\;\text{where}\;\; z_k = \frac{k^{-r}}{\sqrt{|s_k-s_j|}}.
$$
In a similar way, $D^{-1/2}h$ can also be upper bounded component-wise as
$$
\big|[D^{-1/2}h]_k \big| \lesssim \frac{1}{\sqrt{n}}\frac{j^{-r}k^{-r}}{\sqrt{|s_k-s_j|}} = \frac{j^{-r}}{\sqrt{n}}z.
$$

Let $C>0$ be a sufficiently large numerical constant (potentially depending on constants in spectrum decay and spectral gaps). Discuss three cases:
\begin{enumerate}
\item The case of $k\leq j/C$. In this case $|s_k-s_j|\gtrsim s_k \gtrsim k^{-2r}$, and therefore $|z_k|\lesssim O(1)$.
\item The case of $k\geq Cj$. In this case $|s_k-s_j|\gtrsim s_j \gtrsim j^{-2r}$, and therefore $|z_k|\lesssim O(j^rk^{-r})$.
\item The case of $j/C\leq k\leq Cj$. In this case $|s_k-s_j| \gtrsim |k-j| j^{-2r-1}$, and therefore $|z_k|\lesssim O(\sqrt{j/|k-j|})$.
\end{enumerate}
Because $\tilde H=zz^\top/\sqrt{n}$, summarizing the above cases we obtain
\begin{align}
\|\tilde H\|_\op &\lesssim \frac{1}{\sqrt{n}}\|z\|_2^2 = \frac{1}{n}\sum_{k\neq j}|z_k|^2  \lesssim \frac{1}{\sqrt{n}}\left(j + j^{2r}\int_j^{\infty}t^{-2r}\ud t + j\int_1^{Cj} \frac{\ud t}{t}\right) \lesssim \frac{j\ln (j+1)}{\sqrt{n}}.
\end{align}
With $r>1+\sqrt{2}/2=1.207$, we have $J_1=n^{a}$ for some $a<1/2$. Therefore, $\|\tilde H\|_\op = o(1)$.
{\color{blue}[Remark: utilizing the fact that $\mathbb E[\tilde H_{jk}]=0$ we might get tighter control of $\|\tilde H\|_\op$ with high probability. Highly likely $\sqrt{j/n}$. However, it is going to be harder to bound $\tilde H^\ell$ later since we don't have element-wise non-negativity.]}

We can then do the following Taylor expansion using $(I+\tilde H)^{-1} = \sum_{\ell=0}^{\infty} (-1)^\ell \tilde H^\ell$:
\begin{align*}
\|w\|_2 &\leq \sum_{\ell=0}^{\infty} \|D^{-1/2}\tilde H^\ell D^{-1/2}h\|_2 \leq \sum_{\ell=0}^{\infty} \frac{j^{-r}}{\sqrt{n}}\left\|D^{-1/2}\left(\frac{1}{\sqrt{n}}zz^\top\right)^\ell z\right\|_2\\
&\leq \sum_{\ell=0}^{\infty}\frac{j^{-r}}{\sqrt{n}}\left(\frac{j}{\sqrt{n}}\right)^\ell \|D^{-1/2}z\|_2 \lesssim \frac{j^{-r}}{\sqrt{n}}\|D^{-1/2}z\|_2\\
&\lesssim \frac{j^{-r}}{\sqrt{n}}\left(\int_1^j \frac{\ud t}{t^{-2r}} + j^{2r}\int_j^{J_1}\frac{t^{-2r}}{j^{-2r}}\ud t + j\cdot \int_1^{Cj}\frac{1}{tj^{-2r-1}}\frac{\ud t}{t}\right)^{1/2}\\
&\lesssim \frac{j^{-r}}{\sqrt{n}}\left(j^{2r+1} + j^{2r+1} + j^{2r+2}\right)^{1/2} = \frac{j}{\sqrt{n}}.
\end{align*}
With $4r^2-8r-1>0$, implied by $r>2+\sqrt{5}/2=2.118$, $j/\sqrt{n}<1/J_1$ and this proves the desired results.
\end{proof}

\bibliographystyle{plain}
\bibliography{refs}

\end{document}